% Owners: both
\section{Discussion and conclusion}
\label{sec:conclusion}

\paragraph{What limits BFS.} Combining profiles, edge counts and the model gives a
consistent picture. A single core is latency-bound: it cannot keep enough misses in flight.
A laptop with all cores is bandwidth-bound and already beyond its DRAM bound thanks to L3
reuse. A 128-core server is bound by neither: synchronisation at about ten barriers per
search and cross-socket effects dominate. The optimisations confirm this. The significant ones reduce memory traffic or work
(relabelling), synchronisation span (fewer, socket-local threads), or remote memory traffic
(NUMA-aware first touch). Those that target instruction count (compiler flags, vectorisation) or a
minor kernel (prefetching in top-down) were not. This matches the course's message: the
justification from analysis and model predicted which strategies could pay off before they
were implemented.

\paragraph{Threats to validity.}
(i)~The laptop ran under WSL2: no hardware counters, a hidden hybrid topology and high
variance (IQR up to \SI{56}{\percent}). We therefore use medians and treat differences within
the IQR as insignificant. (ii)~The model assumes every byte comes from DRAM; this is why the
laptop beats its bound. A cache-aware refinement would use the L3 bandwidth for the
\SI{4}{MB} \texttt{parent} array. (iii)~Hábrók's counters were multiplexed, and the 128-thread counter runs include warm-up
trials, so we use only single-thread counters quantitatively.
(iv)~Relabelling changes the source vertices. (v)~\texttt{kron20} is small for a 128-core
node. Results on \texttt{kron22}/\texttt{kron24} show the expected higher GTEPS with
growing work per barrier (\SI{6.8}{GTEPS} on \texttt{kron24}).

\paragraph{Future work.} Hardware counters on the laptop through a native Linux boot, which
would also allow a top-down microarchitecture analysis; the Apple Silicon comparison planned
in our work plan; a branch-free bitmap test in \texttt{BUStep} against the
\SI{12.9}{\percent} misprediction rate; a NUMA-aware partitioning of the bottom-up sweep; and
fusing the short parallel regions of a BFS level to remove barriers, analogous to loop
fusion. The last one addresses the dominant cost we identified.
